\SourceLeaf{254}%
\Needspace{7\baselineskip}
\begin{center}
\textsc{Gauss} an \textsc{Schumacher}. Göttingen, 8. Februar 1846.
\end{center}

\mbox{.\ .\ .\ .\ .\ .\ .} Der Unterschied zwischen Rechts und Links lässt sich aber nicht definieren, sondern nur vorzeigen. Insofern verhält es sich damit ähnlich wie mit Süß und Bitter. Omne simile claudicat aber. Letzteres gilt nur für Wesen mit Geschmacksorganen, Ersteres hingegen für alle Geister, die die materielle Welt erfassen können. Zwei solche Geister können sich über Rechts und Links aber nur dann unmittelbar verständigen, wenn ein und dasselbe einzelne materielle Ding eine Brücke zwischen ihnen schlägt. Ich sage unmittelbar, weil sich auch A mit Z verständigen kann, indem zwischen A und B eine materielle Brücke, zwischen B und C eine andere usw. geschlagen wird oder worden sein kann. Welche Bedeutung diese Sache für die Metaphysik hat und dass ich darin die schlagende Widerlegung von \textsc{Kants} Einbildung finde, der Raum sei BLOSS die Form unserer äußeren Anschauung, habe ich knapp in den Göttingischen Gelehrten Anzeigen 1831, S. 635 ausgesprochen. \mbox{.\ .\ .\ .\ .\ .\ .}

\par\bigskip

\Needspace{7\baselineskip}
\begin{center}
[\textsc{Gerling} \emph{an} \textsc{Gauss}. \emph{Marburg, 20. Juni 1846.}]
\end{center}

\{\mbox{.\ .\ .\ .\ .\ .\ .\ .} Zum ersten Mal musste ich eine \textls{rechts}gewundene Schraube für den Druck [der neuen Ausgabe des »Lorenz«] definieren. Ich weiß es nicht besser zu sagen als so: Sie ist so gewunden, dass jeder Punkt der senkrecht stehenden Schraube außerhalb der Achse nur nach unten gelangen kann, wenn er von Ost über Süd nach West geht. --- Gibt es eine bessere Definition, so würden Sie mich sehr verbinden, wenn Sie sie mir mitteilten. \mbox{.\ .\ .\ .\ .\ .\ .}\}

\par\bigskip
